\documentclass[12pt]{article}
\makeatletter
\def\input@path{{../../}}
\makeatother
\usepackage{sbc-template}
\makeatletter
\renewcommand{\inst}[1]{}
\makeatother
\usepackage{graphicx,url}
\usepackage[utf8]{inputenc}
\usepackage[T1]{fontenc}
\usepackage[brazil]{babel}
\usepackage{longtable,booktabs,array}
\providecommand{\tightlist}{%%
  \setlength{\itemsep}{0pt}\setlength{\parskip}{0pt}}
\sloppy
\title{Avaliacao justa em previsao de demanda: baselines, IA classica e Reservoir Computing}
\author{}
\address{}
\begin{document}

\maketitle

\begin{resumo}
Se os artigos anteriores mostraram o problema, a intuicao de RC e o primeiro pipeline reproduzivel, este artigo entrega a espinha dorsal metodologica da serie: um protocolo justo de comparacao entre familias de modelos. A partir do mesmo recorte do Favorita, com a mesma janela de teste e as mesmas metricas, comparamos Seasonal Naive, ETS, Prophet, XGBoost, LSTM, RC e QRC. O benchmark atual coloca ETS na lideranca (\texttt{MAE\ =\ 216.303}, \texttt{RMSE\ =\ 307.638}, \texttt{WAPE\ =\ 0.0999}, \texttt{sMAPE\ =\ 0.1042}), seguido por XGBoost e Prophet. O RC supera a LSTM simples do projeto, mas permanece abaixo dos modelos classicos mais fortes. O QRC, por sua vez, fecha o benchmark com uma implementacao funcional, embora ainda nao competitiva neste recorte. O valor central deste artigo nao esta apenas na tabela final, mas nas regras que tornam a tabela defensavel.
\end{resumo}

\section{1. Introducao}\label{1-introducao}

Comparar modelos temporais parece facil ate o momento em que percebemos quantas comparacoes sao injustas por construcao. Em previsao de demanda, basta um pequeno vazamento de informacao, um split mal definido ou uma metrica escolhida sem criterio para transformar o experimento inteiro em uma ilustraçao enganosa.

Por isso, este artigo assume uma posicao clara: antes de discutir QRC, precisamos construir um benchmark digno de comparacao. Isso significa:

\begin{itemize}
\tightlist
\item
  mesmo recorte de dados;
\item
  mesmo horizonte de teste;
\item
  metricas comuns;
\item
  interpretacao metodologica e de negocio.
\end{itemize}

\section{2. Por que avaliacao em series temporais e delicada}\label{2-por-que-avaliacao-em-series-temporais-e-delicada}

Series temporais nao admitem o mesmo tipo de embaralhamento tipico de problemas tabulares iid. Em forecasting, o tempo nao e apenas mais uma coluna; ele define a direcao causal do problema.

Os erros mais comuns sao:

\begin{itemize}
\tightlist
\item
  usar informacao do futuro na engenharia de atributos;
\item
  normalizar usando o dataset inteiro;
\item
  misturar treino e teste em rolling features mal construidas;
\item
  comparar modelos em horizontes ou variaveis diferentes;
\item
  declarar um vencedor com base em uma unica metrica.
\end{itemize}

Uma boa avaliacao em series temporais precisa ser mais parecida com o ambiente de producao do que com um jogo de leaderboard.

\section{3. Protocolo experimental da serie}\label{3-protocolo-experimental-da-serie}

O protocolo adotado no projeto foi desenhado para ser simples, legivel e reaproveitavel.

\subsection{3.1 Recorte comum}\label{31-recorte-comum}

\begin{itemize}
\tightlist
\item
  dataset: Favorita \texttt{train.csv}
\item
  serie: \texttt{store\_nbr\ =\ 1}, \texttt{family\ =\ BEVERAGES}
\item
  frequencia: diaria
\item
  janela de teste: ultimos 90 dias
\end{itemize}

\subsection{3.2 Regra de split}\label{32-regra-de-split}

Todo treino ocorre antes do bloco final de teste. Nenhum modelo recebe acesso direto a observacoes futuras.

\subsection{3.3 Modelos avaliados}\label{33-modelos-avaliados}

\begin{itemize}
\tightlist
\item
  Seasonal Naive
\item
  ETS
\item
  Prophet
\item
  XGBoost
\item
  LSTM
\item
  RC
\item
  QRC
\end{itemize}

\subsection{3.4 Metricas avaliadas}\label{34-metricas-avaliadas}

\begin{itemize}
\tightlist
\item
  \texttt{MAE}
\item
  \texttt{RMSE}
\item
  \texttt{WAPE}
\item
  \texttt{sMAPE}
\end{itemize}

Essas metricas foram escolhidas porque combinam interpretabilidade absoluta, sensibilidade a erros grandes e leitura relativa ao volume da serie.

\section{4. Benchmark comparativo atual}\label{4-benchmark-comparativo-atual}

Os resultados atuais do projeto sao:

\begin{longtable}[]{@{}llrrrr@{}}
\toprule\noalign{}
Modelo & Familia & MAE & RMSE & WAPE & sMAPE \\
\midrule\noalign{}
\endhead
\bottomrule\noalign{}
\endlastfoot
ETS & baseline estatistico & 216.303 & 307.638 & 0.0999 & 0.1042 \\
XGBoost & ML classico & 256.666 & 332.828 & 0.1185 & 0.1237 \\
Prophet & baseline estatistico & 257.584 & 321.818 & 0.1189 & 0.1319 \\
Seasonal Naive & baseline estatistico & 259.067 & 352.280 & 0.1196 & 0.1361 \\
RC & reservoir computing classico & 324.294 & 423.417 & 0.1498 & 0.1653 \\
LSTM & DL classico & 371.827 & 461.342 & 0.1717 & 0.1822 \\
QRC & reservoir computing quantico & 441.286 & 525.325 & 0.2038 & 0.2269 \\
\end{longtable}

Essa tabela ja faz um trabalho importante para a serie:

\begin{itemize}
\tightlist
\item
  mostra que baselines estatisticos continuam fortissimos;
\item
  impede que modelos mais complexos sejam julgados contra referencias fracas;
\item
  fixa um ponto de comparacao unico para todos os artigos seguintes.
\end{itemize}

\section{5. Leitura dos resultados}\label{5-leitura-dos-resultados}

\subsection{5.1 ETS na lideranca}\label{51-ets-na-lideranca}

O melhor desempenho do ETS neste recorte e pedagogicamente revelador. Ele lembra que, quando a serie tem sazonalidade forte e estrutura relativamente regular, um modelo estatistico bem calibrado pode ser extremamente competitivo.

\subsection{5.2 XGBoost e Prophet como faixa intermediaria forte}\label{52-xgboost-e-prophet-como-faixa-intermediaria-forte}

XGBoost e Prophet aparecem como modelos muito serios para um problema de negocio real. Eles nao eliminam a necessidade de RC, mas elevam bastante o nivel da comparacao.

\subsection{5.3 O que o RC entrega}\label{53-o-que-o-rc-entrega}

O RC nao ganha do topo do benchmark, mas supera a LSTM simples deste projeto. Isso mostra que o paradigma e capaz de sustentar uma representacao temporal competitiva sem o custo de treinar toda a dinamica recorrente.

\subsection{5.4 O que o QRC entrega}\label{54-o-que-o-qrc-entrega}

O QRC ainda nao e competitivo com os melhores modelos classicos neste recorte. Ainda assim, ele entra no benchmark com valor metodologico alto:

\begin{itemize}
\tightlist
\item
  funciona de ponta a ponta;
\item
  respeita o mesmo split temporal;
\item
  usa as mesmas metricas;
\item
  pode ser analisado dentro do mesmo quadro de comparacao.
\end{itemize}

Esse e um ganho importante para uma serie introdutoria. O objetivo nao e inflar o QRC, mas posiciona-lo corretamente.

\section{6. O que significa uma comparacao "justa"}\label{6-o-que-significa-uma-comparacao-justa}

Comparacao justa nao significa exigir que todos os modelos tenham a mesma estrutura interna. Significa exigir que respondam ao mesmo problema sob as mesmas regras.

No contexto desta serie, isso quer dizer:

\begin{itemize}
\tightlist
\item
  mesmo alvo;
\item
  mesmo horizonte de teste;
\item
  mesmas metricas;
\item
  mesma serie;
\item
  nenhuma vantagem informacional indevida.
\end{itemize}

Essa definicao simples evita um numero enorme de confusoes.

\section{7. Checklist de boas praticas}\label{7-checklist-de-boas-praticas}

Para tornar a serie reutilizavel como material didatico, o protocolo do projeto pode ser resumido no seguinte checklist:

\begin{enumerate}
\def\labelenumi{\arabic{enumi}.}
\tightlist
\item
  definir um recorte de negocio claro;
\item
  separar treino e teste por ordem temporal;
\item
  garantir que features usem apenas informacao disponivel no instante;
\item
  comparar com pelo menos um baseline forte;
\item
  reportar mais de uma metrica;
\item
  interpretar resultados junto com custo e complexidade;
\item
  guardar artefatos e configuracoes no repositorio.
\end{enumerate}

Esse checklist e simples o bastante para ser ensinado cedo e forte o bastante para suportar a chegada do QRC.

\section{8. Regras herdadas para a fase QRC}\label{8-regras-herdadas-para-a-fase-qrc}

Quando a serie entra em QRC, algumas tentacoes aparecem: reduzir demais o problema, trocar metrica, mudar serie ou relaxar comparacao. Este artigo fixa a regra contraria. A fase QRC herda o protocolo ja estabelecido.

Assim, o artigo final pode discutir o que muda no reservatorio sem mudar o problema que esta sendo resolvido.

\section{9. Conclusao}\label{9-conclusao}

O maior resultado deste artigo nao e apenas a tabela comparativa. E a instituicao de um padrao metodologico para a obra inteira. Com esse benchmark, o leitor aprende que:

\begin{itemize}
\tightlist
\item
  previsao de demanda exige rigor experimental;
\item
  modelos classicos continuam sendo referencias importantes;
\item
  RC e QRC so fazem sentido quando entram em comparacoes honestas.
\end{itemize}

Esse e o solo sobre o qual a ponte para QRC sera construida.

\section{Entregaveis associados no repositorio}\label{entregaveis-associados-no-repositorio}

\begin{itemize}
\tightlist
\item
  benchmark consolidado: \texttt{code/RESULTS.md}
\item
  baselines: \texttt{code/baselines/}
\item
  modelos classicos: \texttt{code/classical/}
\item
  RC: \texttt{code/rc/}
\item
  QRC: \texttt{code/qrc/}
\end{itemize}

\section{Referencias}\label{referencias}

\begin{enumerate}
\def\labelenumi{\arabic{enumi}.}
\tightlist
\item
  Hyndman, R. J., Athanasopoulos, G. Forecasting: Principles and Practice.
\item
  Taylor, S. J., Letham, B. Forecasting at scale.
\item
  Chen, T., Guestrin, C. XGBoost: A scalable tree boosting system.
\item
  Hochreiter, S., Schmidhuber, J. Long short-term memory.
\end{enumerate}

\section*{Resultados Computacionais Associados}

Os resultados computacionais associados a este artigo estao organizados na subpasta \texttt{computational\_results\_20260402\_222902/}, localizada dentro desta pasta \LaTeX. Esse diretorio contem o arquivo \texttt{REPORT.md}, tabelas em CSV e imagens comparativas apropriadas ao escopo do artigo.

\end{document}
